
\documentclass[11pt]{article}
\usepackage[margin=1in]{geometry}

% Core packages
\usepackage{amsmath,amssymb,amsthm}
\usepackage[spanish]{babel}
\usepackage{tikz-cd}
\usepackage{physics,exercise,cancel}
\usepackage{bm}
\usepackage{dsfont}
\usepackage{multicol}
\usepackage{tikz}
\usetikzlibrary{decorations.pathmorphing, patterns}

% Paragraphs
\setlength{\parindent}{0pt}
\setlength{\parskip}{1\baselineskip}

% Problems
\theoremstyle{definition}
\newtheorem*{problem}{Problema}

% Equation numbering
\numberwithin{equation}{section}

\title{Ayudantía MMF II}
\author{Jovanny Jara}
\date{05 de junio de 2026}

\begin{document}
\maketitle

\section{Ecuación de Laplace en 3D: Un problema de Electrostática}

\begin{Exercise*}[title={Ejemplo 3.9}, name={}, origin={Griffiths, Introduction to Electrodynamics}]
Sea un cascarón esférico de radio $R$ con una densidad de carga superficial $\sigma_0(\theta)$, donde $\theta$ es el ángulo polar. Encuentre el potencial eléctrostático en todo el espacio.
\end{Exercise*}
\begin{proof}[Solución]\renewcommand{\qedsymbol}{}
  Es posible trabajar este problema mediante integración directa, pero dado que tenemos una densidad de carga no uniforme que depende de $\theta$, es conveniente utilizar separación de variables asumiendo simetría azimutal. De esta forma, ya conocemos la solución general a la ecuación de Laplace, la cual es, 
  \begin{equation}
    \Phi(r,\theta) = \sum_{l=0}^{\infty}\left[A_l r^l + B_l r^{-(l+1)}\right]P_l(\cos\theta).
  \end{equation}
  Entonces, solo necesitamos trabajar tanto dentro como fuera del cascarón y aplicar condiciones de borde para determinar los coeficientes $A_l$ y $B_l$. 
  Veamos el caso interior. Como el potencial debe ser finito en el origen, se concluye que $B_l=0$ para todo $l$. Por lo tanto, la solución dentro de la esfera es,
  \begin{equation}
    \boxed{\Phi(r,\theta) = \sum_{l=0}^{\infty}A_l r^l P_l(\cos\theta) \qc (r\leq R)}.
  \end{equation}
  Para el caso exterior, el potencial debe tender a cero a medida que $r$ tiende a infinito, lo que implica que $A_l=0$ para todo $l$. Por lo tanto, la solución fuera de la esfera es,
  \begin{equation}
    \boxed{\Phi(r,\theta) = \sum_{l=0}^{\infty}B_l r^{-(l+1)} P_l(\cos\theta) \qc (r\geq R)}.
  \end{equation}
  Ahora, aplicamos las condiciones de borde. La primera condición es la continuidad del potencial en $r=R$, entonces,
  \begin{equation}
    \sum_{l=0}^{\infty}A_l R^l P_l(\cos\theta) = \sum_{l=0}^{\infty}B_l R^{-(l+1)} P_l(\cos\theta).
  \end{equation}
  Multiplicamos ambos lados por $P_{l'}(\cos\theta)\sin\theta$ e integramos en $\theta$ desde $0$ hasta $\pi$.
  \begin{equation}
    \int_0^{\pi}\sum_{l=0}^{\infty}A_l R^l P_l(\cos\theta)P_{l'}(\cos\theta)\sin\theta \dd{\theta} = \int_0^{\pi}\sum_{l=0}^{\infty}B_l R^{-(l+1)} P_l(\cos\theta)P_{l'}(\cos\theta)\sin\theta \dd{\theta}.
  \end{equation}
  Utilizando la ortogonalidad de los polinomios de Legendre, obtenemos,
  \begin{equation}
    \boxed{B_l = A_l R^{2l+1}}.
  \end{equation}
  La segunda condición de borde se debe a la discontinuidad del campo eléctrico en la superficie del cascarón ya que existe una densidad de carga superficial. Esto se expresa como,
  \begin{equation}
    \eval{\pqty{\pdv{\Phi_{\text{ext}}}{r} - \pdv{\Phi_{\text{int}}}{r}}}_{r=R} = -\frac{\sigma_0(\theta)}{\epsilon_0}.
  \end{equation}
  Entonces, tenemos,
  \begin{equation}
    -\sum_{l=0}^{\infty}(l+1)B_l R^{-(l+2)} P_l(\cos\theta) - \sum_{l=0}^{\infty}l A_l R^{l-1} P_l(\cos\theta) = -\frac{\sigma_0(\theta)}{\epsilon_0}.
  \end{equation}
  Usando la ecuación (1.6), podemos escribir,
  \begin{equation}
    \sum_{l=0}^{\infty}(2l+1)A_l R^{l-1} P_l(\cos\theta) = \frac{\sigma_0(\theta)}{\epsilon_0}.
  \end{equation}
  Luego, para encontrar los coeficientes $A_l$, volvemos a multiplicar ambos lados por $P_{l'}(\cos\theta)\sin\theta$ y aplicamos ortogonalidad de los polinomios, lo que nos da,
  \begin{equation}
    \boxed{A_l = \frac{1}{2\epsilon_0 R^{l-1}}\int_0^{\pi}\sigma_0(\theta)P_l(\cos\theta)\sin\theta \dd{\theta}}.
  \end{equation}
  Así, las ecuaciones (1.2) y (1.3) junto con los coeficientes dados por las ecuaciones (1.6) y (1.10) nos dan la solución completa al problema.
\end{proof}

\section{Ecuación de onda en 1+1 dimensiones: Uno cortito del Arfken}

\begin{Exercise*}[title={Ejercicio 9.6.1}, name={}, origin={Arfken, Weber, Harris, Mathematical Methods for Physicists}]

  Resolver la ecuación de onda unidimensional (estacionaria) y obtener $\phi(x,t)$, si en $t=0$ se tiene que $\phi(x,0) = \sin(x)$ y $\pdv*{\phi(x,0)}{t}= \cos(x)$.

\end{Exercise*}

\begin{proof}[Solución]\renewcommand{\qedsymbol}{}
  La ecuación a resolver es,
  \begin{equation}
    \pdv[2]{\phi}{t} = c^2 \pdv[2]{\phi}{x}.
  \end{equation}
  Sabemos que la solución general a esta ecuación es
  \begin{equation}
    \phi(x,t) = \sum_{n=1}^{\infty}\left[A_n \cos\left(\frac{n\pi c t}{L}\right) + B_n \sin\left(\frac{n\pi c t}{L}\right)\right]\sin\left(\frac{n\pi x}{L}\right),
  \end{equation}
  con 
  \begin{equation}
    A_n = \frac{2}{L}\int_0^L \phi_0(x)\sin\left(\frac{n\pi x}{L}\right) \dd{x} \qc B_n = \frac{2}{n\pi c}\int_0^L \pdv*{\phi(x)}{t}\sin\left(\frac{n\pi x}{L}\right) \dd{x},
  \end{equation}
  donde 
  \begin{equation}
    \phi(x,0) = f(x) \qc \pdv*{\phi(x,0)}{t} = g(x).
  \end{equation}
  En nuestro caso, $f(x) = \sin(x)$ y $g(x) = \cos(x)$. Entonces, calculamos los coeficientes. Para $A_n$, tenemos,
  \begin{equation}
    A_n = \frac{2}{L}\int_0^L \sin(x)\sin\left(\frac{n\pi x}{L}\right) \dd{x}.
  \end{equation}
  Escribimos $k = \frac{n\pi}{L}$ y resolvemos la integral por partes, con $u = \sin(kx)$ y $dv = \sin(x) \dd{x}$, llegando a
  \begin{equation}
    \int_0^L \sin(x)\sin\left(\frac{n\pi x}{L}\right) \dd{x} = k\int_{0}^{L} \cos(x)\cos(kx) \dd{x},
  \end{equation}
  volvemos a integrar por partes, tomando $u = \cos(kx)$ y $dv = \cos(x) \dd{x}$, y finalmente obtenemos,
  \begin{equation}
    A_n = \boxed{\frac{2n\pi(-1)^n\sin(L)}{L^2 - n^2\pi^2}}.
  \end{equation}
  Luego, para $B_n$,
  \begin{equation}
    B_n = \frac{2}{n\pi c}\int_0^L \cos(x)\sin\left(\frac{n\pi x}{L}\right) \dd{x}.
  \end{equation}
  Nuevamente, escribimos $k = \frac{n\pi}{L}$ y resolvemos la integral por partes, con $u = \sin(kx)$ y $dv = \cos(x) \dd{x}$, llegando a
  \begin{equation}
    \int_0^L \cos(x)\sin\left(\frac{n\pi x}{L}\right) \dd{x} = -k\int_{0}^{L} \sin(x)\cos(kx) \dd{x},
  \end{equation}
  volvemos a integrar por partes con $u = \cos(kx)$ y $dv = \sin(x) \dd{x}$ y finalmente obtenemos,
  \begin{equation}
    B_n = \boxed{\frac{2L[(-1)^n\cos(L) - 1]}{c(L^2 - n^2\pi^2)}}.
  \end{equation}
  Con esto, entonces, la solucion completa al problema es
  \begin{equation}
    \phi(x,t) = \sum_{n=1}^{\infty}\sin(\frac{n\pi x}{L})\left[\frac{2n\pi(-1)^n\sin(L)}{L^2 - n^2\pi^2} \cos\left(\frac{n\pi c t}{L}\right) + \frac{2L[(-1)^n\cos(L) - 1]}{c(L^2 - n^2\pi^2)} \sin\left(\frac{n\pi c t}{L}\right)\right].
  \end{equation}
\end{proof}

\end{document}
